\documentclass[12pt]{article}
\usepackage{ctex}
\usepackage{sci-manuscript}
\addbibresource{references.bib}

\ManuscriptTitle{\LaTeX 模板标题}

\ManuscriptAuthors{%
  \AuthorORCID{Lintao Luo}{1,*}{0009-0001-5790-4866}
}

\ManuscriptAffiliations{%
  \Affiliation{1}{giantlinlinlin@gmail.com}%
}

\Keywords{effective population size; Skygrid; trajectory clustering}

% \Correspondence{*Corresponding author: author@example.edu}

\begin{document}

\MakeManuscriptTitle

\section{Materials and Methods}

\subsection{Effective population size trajectories and feature extraction}
Posterior trajectories of the effective population size were taken from
Bayesian Skygrid coalescent analyses \parencite{gill2013} performed in
\software{BEAST} (v1.10) \parencite{suchard2018}, with 200 posterior draws
retained for each of the eight geographic categories after the first
\SI{10}{\percent} of each chain had been discarded as burn-in and the remaining
states had been thinned to at most 4{,}000 draws. All draws were interpolated
onto a common calendar grid running from the present to \SI{20000}{years} before
present in steps of \SI{200}{years}, the upper bound being the oldest time
covered by every input. Each draw was smoothed and differentiated in a single
pass by locally weighted quadratic regression with a tricube kernel and a
bandwidth of \SI{3000}{years} \parencite{cleveland1988}, yielding a smoothed
$\log N_\mathrm{e}$ curve and its rate of change expressed per \SI{1000}{years}.
The time axis was divided into 40 non-overlapping windows of \SI{500}{years},
and eight trajectory features (mean level, net change, mean rate of change,
fluctuation amplitude, and the timing and magnitude of the strongest expansion
and contraction) were computed for every draw within each window; medians,
\SI{95}{\percent} highest posterior density (HPD) intervals and directional
posterior probabilities were then summarised across draws. A grid point was
regarded as reliable when the \SI{95}{\percent} HPD width of $\log N_\mathrm{e}$
did not exceed 2.0 log units, and a category entered a given window only when at
least \SI{80}{\percent} of the grid points of that window were reliable. All
computations were implemented in \software{Python} (v3.10.20) with
\software{NumPy} (v2.2.6) \parencite{harris2020}, \software{SciPy} (v1.15.2)
\parencite{virtanen2020} and \software{pandas} (v2.3.3) \parencite{mckinney2010}.

\subsection{Trajectory clustering, ordination and visualisation}
Within each time window the eight features were standardised across categories
($z$-scores) and clustered by average-linkage hierarchical clustering of
Euclidean distances using \software{SciPy} (v1.15.2) \parencite{virtanen2020};
the number of clusters was selected between $K = 2$ and $K = 8$ by maximising the
average silhouette coefficient \parencite{rousseeuw1987} computed with
\software{scikit-learn} (v1.7.2) \parencite{pedregosa2011}. Cluster labels were
made comparable between consecutive windows by optimal label matching with the
Hungarian algorithm \parencite{kuhn1955}. Uncertainty in the partition was
quantified by 500 posterior resamples in which one curve per category was drawn
at random, the features were recomputed and the clustering repeated with the same
$K$ (random seed 20260828); the proportion of resamples placing two categories in
the same cluster was reported as the co-clustering frequency, and pairs with a
frequency of at least 0.5 were drawn as chords. The scaled feature matrix of each
window was further projected onto its first two principal components
\parencite{pedregosa2011}, and two composite scores (a level score from mean
level and amplitude, and a change score from net change, rate and the extreme
expansion and contraction rates) were defined as the first principal component of
the corresponding feature subsets and standardised within each \SI{500}{years}
bin. Figures were produced in \software{R} (v4.6.0) \parencite{rcoreteam2026}
with \software{tidyplots} (v0.3.1) \parencite{engler2025}, \software{ggplot2}
(v4.0.3) \parencite{wickham2016}, \software{ggrepel} (v0.9.8)
\parencite{slowikowski2025}, \software{pheatmap} (v1.0.13) \parencite{kolde2025},
and \software{igraph} (v2.3.3) \parencite{csardi2006} together with
\software{ggraph} (v2.2.2) \parencite{pedersen2024}.

\printbibliography[title={References}]

\end{document}
